\documentclass[11pt]{article}

\usepackage[T1]{fontenc}
\usepackage[utf8]{inputenc}
\usepackage{amsmath,amssymb,amsthm,mathtools}
\usepackage[margin=1in]{geometry}
\usepackage[hidelinks]{hyperref}
\usepackage{booktabs}
\usepackage{xcolor}
\usepackage{listings}

\lstset{
  basicstyle=\ttfamily\small,
  breaklines=true,
  columns=fullflexible,
  frame=single,
  keepspaces=true,
  keywordstyle=\bfseries,
  commentstyle=\itshape\color{gray}
}

\newtheorem{definition}{Definition}
\newtheorem{proposition}{Proposition}
\newtheorem{claim}{Claim}

\newcommand{\CDC}{BiDi Coherence-Delta Calculus}
\newcommand{\C}{\mathcal{C}}
\newcommand{\eps}{\varepsilon}
\newcommand{\flow}[1]{\xrightarrow{#1}_{d}}
\newcommand{\commit}{\xrightarrow{\beta}}
\newcommand{\parcomp}{\mathbin{\|}}
\newcommand{\rot}{\circlearrowright}
\newcommand{\botv}{\bot}
\newcommand{\bidi}{\mathrm{bidi}\gamma\Delta}
\newcommand{\nest}[1]{\mathopen{[\![}#1\mathclose{]\!]}}

\title{\CDC:\\A Compact Executable Calculus for Hybrid Continuous-Discrete Coherence}
\author{Edward Ellis\\Independent Researcher}
\date{June 18, 2026}

\begin{document}
\maketitle

\begin{abstract}
We introduce \CDC, a small formal substrate for systems that combine continuous
dynamics, event-triggered balanced-ternary commitments, delayed angular coupling, local
invariants, nested reference frames, and belief/error minimization. The calculus
uses five canonical primitives---cells, channels, modules, fields, and
commits---plus a bidirectional coherence-delta operator for coupling nested
scales and path endpoints. Its current executable artifact includes a
transitional Python reducer, a native source language (\texttt{.cdc}),
executable law witnesses, and
capability witnesses for hybrid flow, off-grid commits, delay lines,
angle-biased relation channels, multirate nesting, predictive belief updates,
local plasticity, local trace-order windows, balanced-ternary trace/window measurement, and a two-counter
universality construction. The
current artifact is not presented as a completed theorem-prover development;
rather, it is an executable calculus whose algebraic and reduction properties
are stated compactly and checked by a reproducible witness suite.
\end{abstract}

\section{Motivation}

Software systems that model embodied agents, mixed-reality worlds, adaptive
interfaces, simulations, and reactive infrastructure often need several regimes
at once: continuous flow, discrete state transitions, delayed influence,
invariants, nested scale, belief updates, and local policy gates. These regimes
are usually split across unrelated tools: dynamical systems, event buses, state
machines, workflow runtimes, neural models, policy engines, and dataflow graphs.

\CDC\ is an attempt to compress these regimes into one executable formal object.
The goal is not to replace every specialized runtime. The goal is to give hybrid
systems one shared semantic spine:
\[
\begin{aligned}
\text{carrier algebra} &\to \text{terms} \to \text{congruence}\\
&\to \text{flow/commit reduction} \to \text{operator algebra}\\
&\to \text{executable witnesses}.
\end{aligned}
\]

Related trace-order programs model observation through observer windows over
Markov kernels and treat observer time as local ordered experience rather than a
detached global clock \cite{hoffman2024traces,hoffman2014time}. \CDC\ does not
adopt that ontology. It only preserves the engineering constraint that
observation is a bounded role over a trace, not a privileged external primitive.

\paragraph{Expository order.}
The presentation keeps the formal object, source notation, and reference reducer
close to one another. Definitions are introduced where their executable witnesses
first become visible, so the reader can follow the calculus as both a mathematical
specification and a running artifact.

\section{Carrier Algebra}

\begin{definition}[Coherence algebra]
A coherence algebra is a structure
\[
\C = (C,\oplus,\botv,\rot,|\cdot|,\kappa)
\]
where $(C,\oplus,\botv)$ is a commutative monoid, $\rot_\phi:C\to C$ is a circle
action by $\oplus$-automorphisms, $|\cdot|:C\to\mathbb{R}_{\geq 0}$ is an
amplitude seminorm, and $\kappa:C\to[-1,1]$ is a commitment coordinate.
\end{definition}

The standard model used by the reference reducer is $C=\mathbb{R}^2 \simeq
\mathbb{C}$, $\oplus$ is vector addition, $\botv=0$, $\rot_\phi$ is phase
rotation, and $\kappa(x)$ is the cosine coordinate of $x$ when $x\ne 0$.
For a fixed deadband $\delta\in(0,1)$, define the runtime trit
\[
\tau(x)=
\begin{cases}
+ & \kappa(x)>\delta,\\
- & \kappa(x)<-\delta,\\
0 & \text{otherwise,}
\end{cases}
\]
and the boundary openness
\[
\chi(x)=\max(0,1-|\kappa(x)|/\delta).
\]

\section{Terms}

Fields are generated by five constructors:
\[
F,G ::= 0
  \mid m[\bar v;\bar b]
  \mid F\parcomp G
  \mid p \triangleright \langle w,t,\alpha,L\rangle p'
  \mid m\nest{F} .
\]

Here $m[\bar v;\bar b]$ is a module containing a cell vector $\bar v\in C^n$ and
belief vector $\bar b\in\mathbb{R}^n$; $F\parcomp G$ is parallel composition;
$p\triangleright\langle w,t,\alpha,L\rangle p'$ is a delayed weighted relation
from source path $p$ to destination path $p'$, with angular phase bias $\alpha$
and optional target-line projection $L$; and $m\nest{F}$ is nested field
structure.

Structural congruence makes $(F,\parcomp,0)$ a commutative monoid and treats
channels as an order-free multiset. This is the compositional backbone: parallel
order is not semantic order.

\section{Hybrid Reduction}

Computation is an interleaving of continuous flow and guarded discrete commits.

\paragraph{Flow.}
$F\xrightarrow{d}F'$ evolves phases, amplitudes, beliefs, and plastic weights for
real duration $d\geq 0$ under an afferent-driven vector field. Each incoming
relation reads its source path at delay $t$, rotates that phasor by $\alpha$ into
the destination frame, optionally projects it to target lines $L$, and gates it
by source and target openness before superposition. In the reference reducer,
phase uses RK4 integration while slow variables use forward Euler. This is an
implementation of the stated semantics, not the definition of the calculus.

\paragraph{Commit.}
$F\commit F'$ applies an evented guarded update to a module:
\begin{enumerate}
\item phase cells snap fractionally toward committed poles;
\item a nonnegative prefix-walk invariant is enforced over trits;
\item belief steps toward local evidence;
\item the update is rejected if local free energy increases;
\item committed nonzero trits latch persistent poles.
\end{enumerate}

The local objective is
\[
\Phi_m =
\sum_i \left[
\frac{1}{2}\pi(\operatorname{Re}\hat A_i-b_i)^2
+\frac{1}{2}(b_i-b_i^0)^2
-|\hat A_i|\cos(\arg\hat A_i-\theta_i)
\right].
\]

\paragraph{Trace windows and measurement.}
A window is a derived path/line/angle-bounded projection over a field. A trace
records phase-time and event-time through that window without changing dynamics.
A committing measurement is a windowed commit record: it uses the existing
guarded commit, reports a balanced-ternary outcome vector in
$\{-1,0,+1\}^n$, and must not increase $\Phi_m$. Thus observer, participant,
and measurement interface are roles over windows, not new primitive entities.
Trace order is local to the window: smooth phase motion can accumulate with zero
commit events, and different windows can see different event densities over one
field.

\paragraph{Nested reference frames.}
For $m\nest{F}$, \CDC\ couples parent and child by bidirectional
coherence-delta, written $\bidi$: parent context is carried downward and child
coherence contributes afferent evidence upward. Operationally these are the
neutral $\alpha=0$ aggregate relations of the same path-channel operator.
Nonzero $\alpha$ represents angled reference-frame transfer; projected $L$ names
the destination dimensions participating in the relation. Distinct nesting paths
are distinct reference frames.

\section{Operator Algebra}

The invariant registry has thirteen keys. Four pin the finite bridge/carrier,
existence, and trace-order layers, and four are algebraic invariants over the
standard carrier:

\begin{center}
\begin{tabular}{llll}
\toprule
Operator & Symbol & Meaning & Witnessed laws\\
\midrule
balanced carrier & $\{-1,0,+1\}$ & equilibrium-centered commits & finite codomain\\
dyadic/triadic bridge & $2^6=4^3$ & 64-state codebook & finite bijection\\
existence viability & $E_0$ & bounded coherent continuity & agency spectrum\\
trace-order locality & $T_6$ & local trace time & no global event tick\\
gate & $\odot$ & phase composition & abelian group on $\mathbb{T}^n$\\
interfere & $\oplus$ & superposition & commutative monoid\\
rotation & $\rot_\phi$ & global phase action & linear over $\oplus$\\
core-fold & $\partial$ & latent projection & linear and equivariant\\
\bottomrule
\end{tabular}
\end{center}

\section{Executable Reduction Witnesses}

The remaining five registry keys are reduction and metatheorem invariants. They
are intentionally phrased as executable witnesses rather than completed
proof-assistant theorems.

\begin{claim}[Preservation]
Every guarded commit produces an admissible committed module: the prefix walk of
runtime trits never goes negative.
\end{claim}

The witness now includes both randomized commits and exact enumeration of all
$3^6$ committed balanced-trit walks.

\begin{claim}[Commit soundness]
The commit relation never increases local free energy because the guard holds
instead of accepting a free-energy-increasing update.
\end{claim}

\begin{claim}[Local confluence]
Commits on modules with disjoint channel footprints commute.
\end{claim}

\begin{claim}[Flow additivity]
For off-grid split durations in the reference reducer,
$\xrightarrow{d_1};\xrightarrow{d_2}$ agrees with $\xrightarrow{d_1+d_2}$
within numerical realization tolerance.
\end{claim}

\begin{claim}[Trace-order locality]
Smooth phase-time can flow through a bounded window with zero commit events, and
event order remains local to the observing window rather than a universal tick.
\end{claim}

\begin{claim}[Normal forms]
Committed irreducible localized modules correspond to closed nonnegative walks.
For $n=6$, the saturated localized count is Catalan $C_3=5$ and the crossing-
permitting localized count is $51$.
\end{claim}

Together, the four bridge/carrier/existence/trace invariants, four algebraic
invariants, and five reduction invariants make the 13-key semantic registry
exercised by the law checker. The witness suite reports 22/22 law, metatheorem,
viability, trace-order, and bridge checks, 13/13 invariant keys exercised, 32/32 native
\texttt{.cdc} expectations, 5/5
relational phase-channel witnesses, 12/12 balanced-ternary trace/window witnesses, and
24/24 capability witnesses.

\section{Native Source Language}

\CDC\ includes a compact source language, \texttt{.cdc}. A source file declares a
field, modules, channels, guards, reductions, counter programs, and assertions.

\begin{lstlisting}[language={},caption={A small .cdc program.}]
deadband 0.5
field demo dt=0.02 gain=1.4
  module A theta 0 0.3 0.6 0.9 1.2 1.5 omega 1.0
  module B trits + o - + o -
  channel A -> B delay=0.2 weight=1.0 angle=pi/4 lines=0,2,4
  guard B crossing 0
  flow 3.0
  commit B
  expect admissible B
end
\end{lstlisting}

The bridge \texttt{cdc\_boot.py} maps \texttt{.cdc} source onto the reference
reducer. It is an execution bridge, not a separate semantics.

\section{Universality and Boundaries}

The artifact includes a two-counter construction expressed through generic
instruction modules and commits. This witnesses standard register-machine
expressivity in the Minsky sense. It should be read as an expressivity result,
not as a claim that the current reducer is a production compiler.

Likewise, the neural-dynamics witnesses demonstrate native support for recurrent
continuous dynamics, predictive belief updates, action-style correction, Hebbian
weight adaptation, and nested multiscale feedback. They are not claims of
state-of-the-art biological modeling or training efficiency.

\section{Implementation Artifact}

The repository contains:
\begin{itemize}
\item \texttt{BIDI\_CALCULUS\_CORE.md}: formal core;
\item \texttt{FORMAL\_SEMANTIC\_SPINE.md}: AST, runtime-state, reduction, and invariant-spine contract;
\item \texttt{NATIVE\_SELF\_HOSTING\_MANDATE.md}: host-debt burn-down path toward native \texttt{.cdc} self-hosting;
\item \texttt{VERIFICATION\_OBLIGATION\_MATRIX.md}: claim-to-witness-to-proof tracking;
\item \texttt{CDC\_LANGUAGE.md}: source language;
\item \texttt{kernel.cdc}: native self-hosting contract;
\item \texttt{bidi\_calculus.py}: transitional reference reducer;
\item \texttt{cdc\_semantics.py}: declarative AST/state/invariant records;
\item \texttt{cdc\_boot.py}: \texttt{.cdc} bridge;
\item \texttt{calculus\_laws.py}: law/metatheorem witnesses;
\item \texttt{acceptance.py}: capability witnesses;
\item \texttt{relation\_witness.py}: angular/path relation witnesses;
\item \texttt{trace\_window\_witness.py}: balanced-ternary observer-window witnesses;
\item \texttt{system.cdc}, \texttt{laws.cdc}, and \texttt{relations.cdc}: native witness programs.
\end{itemize}

\section{Semantic Spine}

The artifact is organized around one semantic spine rather than parallel
descriptions. \texttt{.cdc} source is parsed into a source AST; execution is
understood against a runtime state tuple; reductions are classified as flow,
commit, or nesting steps; and law checks discharge named invariant records. The
reference reducer remains the current executable realization, while the source
language, paper, and witness suite are kept close enough to be audited as
projections of the same object. The self-hosting target is stronger: move the
parser state, reducer state, witness state, and trace/window state into native
\texttt{.cdc} terms until the host is only a replaceable loader.

This makes the theorem-prover lift narrow and explicit. The finite discrete layer
is the natural first target: trit quantization, prefix-walk admissibility, commit
barrier preservation, and localized normal forms. The carrier algebra and
continuous-flow assumptions can then be formalized around the same invariant
table.

\section{Conclusion}

\CDC\ is a compact executable calculus for hybrid continuous-discrete coherence.
Its contribution is a unified semantic vocabulary and working witness suite for
systems that must flow continuously, commit discretely, preserve local
invariants, and couple nested reference frames bidirectionally.

\begin{thebibliography}{11}

\bibitem{henzinger1996}
T. A. Henzinger.
\newblock The theory of hybrid automata.
\newblock In \emph{Proceedings of the 11th Annual IEEE Symposium on Logic in
Computer Science}, 1996.

\bibitem{milner1989}
R. Milner.
\newblock \emph{Communication and Concurrency}.
\newblock Prentice-Hall, 1989.

\bibitem{minsky1967}
M. L. Minsky.
\newblock \emph{Computation: Finite and Infinite Machines}.
\newblock Prentice-Hall, 1967.

\bibitem{friston2010}
K. Friston.
\newblock The free-energy principle: a unified brain theory?
\newblock \emph{Nature Reviews Neuroscience}, 11:127--138, 2010.

\bibitem{knuth1984lp}
D. E. Knuth.
\newblock Literate programming.
\newblock \emph{The Computer Journal}, 27(2):97--111, 1984.

\bibitem{knuth1984texbook}
D. E. Knuth.
\newblock \emph{The TeXbook}.
\newblock Addison-Wesley, 1984.

\bibitem{hoffman2024traces}
D. D. Hoffman, C. Prakash, and S. Chattopadhyay.
\newblock Traces of consciousness.
\newblock Preprints.org, 2024. DOI: 10.20944/preprints202410.1305.v1.

\bibitem{hoffman2014time}
D. D. Hoffman.
\newblock The origin of time in conscious agents.
\newblock Manuscript, University of California, Irvine, 2014.

\end{thebibliography}

\end{document}
